\chapter{系统架构与 GPU 索引构建}
\label{ch:architecture}
本章回答 RQ1：如何从高维向量构造可搜索的 GPU 常驻 JHQ 索引。变换与主码／残差表示继承自 JHQ\cite{jhq}。系统问题在于训练这一表示、编码完整 FP32 数据超出可用设备空间的数据库，并为查询路径提供一致的列表顺序数组。本章依次讨论训练、放置、分配与验证；查询 kernel 细节见\Cref{ch:gpu-execution}，实测比较见\Cref{sec:7-7}。

\section{端到端建库与搜索边界}
\label{sec:4-1}
索引有三个生命周期阶段：训练生成可复用模型状态，添加编码并排列数据库行，搜索使用完成后的数组。实验必须区分这些阶段。训练状态缓存避免重新学习模型，但仍需要数据库编码与列表构建。冷训练加添加，与复用状态后的查询，测量的是不同操作。

\begin{figure}[tbp]
\centering
% Audited lifecycle; no proposed storage or distributed subsystem.
\begin{tikzpicture}[
 font=\small\linespread{1}\selectfont,
 box/.style={draw=ThesisNavy!65,rounded corners=2pt,align=center,
   text width=3.65cm,minimum height=1.9cm,inner sep=5pt},
 gpu/.style={box,fill=ThesisNavy!7},
 arr/.style={-{Latex[length=2mm]},draw=ThesisNavy!80,thick}]
\node[box] (host) at (0,0) {Host input vectors\\Training samples\\Database batches};
\node[box] (model) at (4.8,0) {Host model\\QR rotation $\Pi$\\Primary codebook};
\node[box] (cache) at (9.6,0) {Host model cache\\$\Pi$ and codebooks\\IVF centroids};
\node[gpu] (train) at (0,-3.6) {GPU training\\Rotate samples\\Fit IVF centroids\\Fit residual model};
\node[gpu] (add) at (4.8,-3.6) {GPU bulk build\\1 or 2 passes\\Assign / encode\\Reorder / pack\\Floats $\leq768$ MiB};
\node[gpu] (index) at (9.6,-3.6) {GPU-resident index\\Offsets and IDs\\Primary / residual codes\\Corrections};
\node[box] (query) at (0,-7.2) {Host query batch\\Pinned staging\\Transfer to GPU};
\node[gpu] (search) at (4.8,-7.2) {GPU query pipeline\\Rotate; probe; scan\\Refine; final top-$k$\\Workspace / graph};
\node[box] (result) at (9.6,-7.2) {Host query results\\IDs and scores\\Transfer from GPU};
\draw[arr] (host) -- node[left,font=\footnotesize] {samples} (train);
\draw[arr] (model.south west) -- (train.north east);
\draw[arr] (train) -- node[above,font=\footnotesize] {model} (add);
\draw[arr] (host.west) -- ++(-.5,0) |- (4.8,-5.15)
 node[pos=.8,below,font=\footnotesize] {database batches} -- (add.south);
\draw[arr] (model) -- (add);
\draw[arr] (add) -- (index);
\draw[<->,dashed,draw=ThesisNavy!80,thick]
 (train.north east) -- node[pos=.8,above,sloped,font=\footnotesize,fill=white,inner sep=1pt]
 {save / load model} (cache.south west);
\draw[arr] (query) -- (search);
\draw[arr] (search) -- (result);
\draw[arr] (index.south) -- ++(0,-.6) -| (search.north);
\end{tikzpicture}

\caption[构建与查询架构]{已实现的建库与查询边界。主机负责采样、正交矩阵构造及解析主码码本生成；评测配置在 GPU 上拟合 IVF 与残差模型。缓存保存训练状态，完成后的数据库表示常驻设备。虚线表示训练状态复用。暂存标注仅约束源码默认的原始／旋转浮点缓冲区，其他流水线 scratch 另计。}
\label{fig:4-architecture}
\end{figure}

主机持有输入向量并提供训练样本，构造方形变换和解析主码码本。GPU 矩阵乘法应用变换；设备训练路径拟合粗粒度 IVF 中心及标量残差码本。添加阶段传输有界批次，编码两级表示，确定 IVF 成员关系，再写成统一的列表顺序表示。因此，主机编排与设备计算共同组成 GPU 常驻搜索系统。

模型状态包含变换、主码码本、残差码本和 IVF 中心。数据库状态进一步包含列表偏移、原始行 ID、主码、残差码及修正项。搜索工作区的逻辑寿命更短：旋转查询、中心分数、探针偏移、LUT 和保留候选缓冲区在下一批被覆盖，但分配的内存保留复用。捕获的 CUDA graph 属于执行状态，引用这些现有数组。

API 通过一次成功的 \texttt{add} 构造固定快照，并拒绝后续 add。搜索该快照时物理位置保持固定。数据变更通过重建替代快照处理，形成简单的批量重建生命周期；可变插入／删除及快照的事务发布还需要额外索引管理。

搜索接收主机查询并返回主机ID及计算分数。分数中的查询常数及其对范围阈值的影响
见\cref{eq:5-refined-constant}，实测主机到主机调用边界见\cref{sec:3-3}。


该设计从主机内存分批传输构建输入，同时将压缩数据库保留在 GPU。它没有提供依赖磁盘的搜索路径，也不意味着任意大的压缩索引都能放入单个设备。


\section{训练、旋转与码本生成}
\label{sec:4-2}
\subsection{变换与采样}

实现在主机生成带 seed 的高斯矩阵，通过 LAPACK QR 构造方形正交矩阵 $\Pi$。对行优先训练数据 $X_t$，设备矩阵乘法计算
\begin{equation}
 Y_t=X_t\Pi^{\mathsf T},\qquad y=\Pi x,\qquad
 \|\Pi(q-x)\|_2^2=\|q-x\|_2^2
 \quad\text{when }\Pi^{\mathsf T}\Pi=I.
 \label{eq:4-transform}
\end{equation}
输出保留维数 $d$。这是正交坐标变换；QR 构造及设备乘法存在有限精度误差，其测量属于\cref{sec:4-6}的 C0。

测试程序按固定步长选择数据库行。主码／IVF样本为Vogue 159,744行，arXiv及两个
OpenAI各319,488行，BGE与Stella各1,277,952行，均为$39n_{\rm list}$。
这与FAISS默认的每质心39训练点警告阈值一致\cite{faiss_clustering}，是采样规则，
不是八轮IVF拟合的统计保证。残差训练另用100,000行等距样本，成本不随较大粗训练样本增长。

历史受控扫描固定主模型和 IVF 中心，仅改变 Vogue 与 OpenAI-3072 的残差训练样本量，覆盖 $B_r\in\{4,8\}$ 和 32/128/256 探针。12 个配置中，100K 相对全量的召回率差值最大为 $0.0018$：11 项较低、1 项相等。增加到中间样本量没有带来单调提升，残差码本身也发生较大变化。这支持有界训练成本的取舍，同时保留可测量的质量损失。该扫描使用 25 轮残差 Lloyd，不能替代之后六个数据集的评测。历史诊断及输入哈希见 \path{data/ch4/review-audit.json}。

\subsection{为何保留解析主码模型}

主码构造器由样本估计共同尺度：
\begin{equation}
 \sigma^2=\frac{1}{n_t d}\sum_{i=0}^{n_t-1}\|x_i\|_2^2,
 \qquad D_s=d/M,\qquad K=2^B.
 \label{eq:4-scale}
\end{equation}
在 $B=D_s=8$ 时，笛卡尔码本使用 $\{-a,+a\}$ 两级，$a=0.674\sigma$；每字节表示八个符号位。这保留 JHQ 的解析构造\cite{jhq}，支持逐坐标 $O(d)$ 编码，以及\cref{sec:5-2}的小 LUT 精确因子化。对每子空间 256 个码字自由执行 Lloyd 优化，通常会破坏这种可分离性。

固定符号重建 $\widehat p=a s$，$s_j\in\{-1,+1\}$ 时，
\begin{equation}
 \|q'-a s\|^2=\|q'\|^2+da^2-2a\langle q',s\rangle.
 \label{eq:4-sign-ranking}
\end{equation}
当 $a>0$，这些固定符号码的实数算术主码排序与尺度无关。这一代数性质无需高斯坐标假设。但尺度仍改变残差 $r=\Pi x-\widehat p$、拟合残差码本和精排顺序；$a=0$ 是退化情况，浮点并列仍须处理。

早期前沿诊断说明了取舍。在 Vogue、$M=96$、$B_r=4$、128 探针下，早期构建与解析主码构建的 primary-only recall 分别为 0.6618 和 0.6342，而完整系统为 0.9452 和 0.9409。残差精排缩小了观测差距。这是包含残差训练变化的历史跨版本比较，不能作为仅替换主码码本的单因素因果估计。本设计选择的是 GPU 执行成本较低的结构化表示；第七章评价其完整质量与成本前沿。

\subsection{粗分配的过滤与精算}

IVF 中心从变换样本的等距行初始化，执行配置的八轮更新。分配按 $s_\ell=\|\mu_\ell\|^2-2\langle y,\mu_\ell\rangle$ 排序。源码提供 FP32 与加速分配路径；满足通常的维数约束时，默认用 int8 乘积过滤，再用 FP32 重算有歧义的选择。

\begin{lemma}[具有认证分数区间的分配]
\label{lem:4-assignment-bound}
令 $\widetilde y,\widetilde\mu_\ell$ 为反量化近似，$e=y-\widetilde y$，$f_\ell=\mu_\ell-\widetilde\mu_\ell$。保持中心范数精确时，近似分数
$\widetilde s_\ell=\|\mu_\ell\|^2-2\langle\widetilde y,\widetilde\mu_\ell\rangle$
的实数算术误差满足
\begin{equation}
 E_\ell=2\bigl(\|y\|\|f_\ell\|+\|e\|\|\widetilde\mu_\ell\|\bigr).
 \label{eq:4-assignment-bound}
\end{equation}
若每个目标分数均处于有效区间 $[L_\ell,U_\ell]$，仅丢弃 $L_\ell>\min_j U_j$ 的中心，并精确处理剩余分数，则保留最小值及约定的并列规则。
\end{lemma}
\begin{proof}
将内积误差展开为 $\langle y,f_\ell\rangle+\langle e,\widetilde\mu_\ell\rangle$，应用 Cauchy--Schwarz。被丢弃分数大于某个保留分数的上界，因此不能取得最小值；保留区间相等情形则不会丢失可能的并列者。
\end{proof}

实现用保守范数界与 FP32 余量常数实现这一过滤与精算模式。引理说明有效界为何允许精确排除；若要求与每种 FP32 累加顺序逐位一致，还需要覆盖相应算术误差的界。已归档诊断对比 FP16 路径，在 Vogue、BGE、Stella 尾批次分别记录 0/14,838、6/31,748、3/16,359 个分配差异。这是有限实现证据，不是完整 FP32 参考认证；冷建日志也未追踪实际执行的分配精度。

\subsection{残差拟合与迭代预算}

IVF 拟合先于残差拟合。后者用固定主模型编码样本，形成 $r=\Pi x-\widehat p(x)$，独立于分配的 IVF 中心。每个子空间拟合一个含 $K_r=2^{B_r}$ 项的有序标量码本，由其 $D_s$ 个坐标共享。GPU 估计器排序残差值，再执行一维 Lloyd 更新；评测配置的上限为 2,000 轮。


历史100K样本诊断支持质量收益递减下的固定预算选择：2,000 相对 25 轮提升 $0.0063$--$0.0076$ recall。后续变化可能超过 BGE 与 Stella 已报告的两次运行波动，因此该诊断未确立统计上严格的平台。诊断还观察到稀疏尾部 bin 持续变化；这是对所实现浮点更新的观察，不是精确 Lloyd 下降定理失效。包括 Stella 在内的冷日志达到上限，因此将训练描述为有预算的迭代拟合。早期协议在上限 20,000 时，Stella 全量与 100K 样本的结果分别为 0.9921 和 0.9920。

\subsection{构建工作量与观测不均衡}

稠密QR一次成本为$O(d^3)$，旋转为$O(n_t d^2)$；直接粗分配每轮为
$O(n_t n_{\rm list}d)$。解析主码构造写入$O(MKD_s)$个值，残差提取访问
$O(n_r d)$个坐标，之后排序并进行标量更新。这些表示工作量，并非延迟预测。
同步阶段诊断将BGE／Stella约83\%训练时间归于IVF k-means（约2.97\,s）。
它使用一个添加缓冲区并强制同步，独立于\cref{sec:7-7}中的正常重叠计时。


列表不均衡进入\cref{sec:5-5}的扫描规则；八轮IVF、等距种子和有限采样
各自对该不均衡的贡献尚未单独隔离。


\paragraph{被测实现的算术精度。}
评测实现v69使用FP32矩阵存储，但在$d\geq1{,}536$时启用cuBLAS TF32 tensor-core
算术：编译期faithful分支将覆盖值置空，随后落入按维度选择的默认分支。
这影响两个OpenAI数据集在训练、add、查询时的旋转及质心GEMM，环境变量无法覆盖
该编译分支。因此FP32存储不代表完整FP32乘法。FP32变体v72修正了这一开关；
敏感性实验复用同一TF32训练缓存，只改变add／查询算术，没有以FP32重新训练模型
（\cref{sec:7-precision}）。


\section{IVF 分配、ID 与物理内存布局}
\label{sec:4-3}
\subsection{从数据库 ID 到列表顺序位置}

输入行 ID 为 $i\in\{0,\ldots,N-1\}$。粗分配根据计算所得中心分数，将变换后的向量放入某个列表。数学分配目标是
\begin{equation}
 a(i)=\operatorname*{arg\,min}_{0\leq\ell<n_{\rm list}}
       \|\Pi x_i-\mu_\ell\|_2^2.
 \label{eq:4-assignment}
\end{equation}
具体实现使用选定的数值分配路径。主码与残差表示编码不减去 $\mu_{a(i)}$；列表成员关系独立于重建，限制之后的搜索集合。

构建器将分配结果与输入 ID 序列共同排序。令 $\rho(p)$ 为排序后物理位置 $p$ 对应的输入 ID，$o_\ell$ 为列表 $\ell$ 的起点。完成后的数组满足
\begin{equation}
 0=o_0\leq o_1\leq\cdots\leq o_{n_{\rm list}}=N,
 \qquad p\in[o_\ell,o_{\ell+1})\Longrightarrow a(\rho(p))=\ell.
 \label{eq:4-offsets}
\end{equation}
偏移通过已排序分配值上的 lower-bound 查找获得。相邻偏移相等表示空列表。按分配结果排序不保证不同独立构建中，同一列表内的行具有稳定顺序。位置在完成后的快照中稳定，其含义是全部关联数组共同描述的行 $\rho(p)$。

\begin{figure}[tbp]
\centering
% Illustrative common-position join; no experiment values.
\begin{tikzpicture}[font=\small\linespread{1}\selectfont,
 cell/.style={draw=ThesisNavy!55,minimum width=1.05cm,minimum height=.65cm,
  inner sep=2pt,align=center},
 lab/.style={anchor=east,align=right,font=\footnotesize}]
\node[anchor=west,font=\small] at (-2.2,1.9)
 {Offsets: $[0,3,3,8]$; list 1 is empty};
\draw[decorate,decoration={brace,amplitude=4pt}] (-.52,1.15) -- (2.62,1.15)
 node[midway,above=5pt,font=\footnotesize] {list 0};
\draw[decorate,decoration={brace,amplitude=4pt}] (2.63,1.15) -- (7.87,1.15)
 node[midway,above=5pt,font=\footnotesize] {list 2};
\node[lab] at (-.65,.6) {Position $p$};
\node[lab] at (-.65,-.25) {ID $\rho(p)$};
\node[lab] at (-.65,-1.15) {Primary word\\group $j=0$};
\node[lab] at (-.65,-2.05) {Residual row};
\node[lab] at (-.65,-2.95) {Correction};
\foreach \p/\id in {0/5,1/1,2/7,3/0,4/6,5/3,6/4,7/2} {
 \node at (1.05*\p,.6) {$\p$};
 \node[cell,fill=ThesisNavy!5] at (1.05*\p,-.25) {$\id$};
 \node[cell] at (1.05*\p,-1.15) {$w_{\p}$};
 \node[cell] at (1.05*\p,-2.05) {$r_{\p}$};
 \node[cell] at (1.05*\p,-2.95) {$\gamma_{\p}$};
}
\draw[ThesisNavy,thick,rounded corners=2pt]
 (2.61,.13) rectangle (3.69,-3.33);
\node[align=center,text width=11.4cm,font=\footnotesize] at (2.8,-4.15)
 {At $p=3$, all fields describe input ID 0.\\
  $w_3=\mathsf w[0,3]$ contains $Z[0,3],Z[1,3],Z[2,3],Z[3,3]$;\\
  $r_3$ denotes a row of $b_r$ bytes, not one residual byte.};
\end{tikzpicture}

\caption[按列表排列的物理索引布局]{同一物理位置连接 ID、打包主码 word、残差行与修正项。示意的八行索引偏移为 $[0,3,3,8]$，其中有一个空列表。主码 word 打包同一向量的四个子空间。数值与列表大小均为示意。}
\label{fig:4-layout}
\end{figure}

常驻 ID 数组实现 $p\mapsto\rho(p)$。逆映射 $\rho^{-1}(i)$ 在数学上存在，但搜索无需再常驻一份逆数组。两遍构建器临时构造逆映射，将输入批次散写到最终位置。查询结果把物理位置翻译回输入 ID；第五章在主码并列时使用位置，在最终并列时使用数据库 ID。

\subsection{精确地址与表示宽度}

令 $Z[m,p]$ 表示物理位置 $p$ 在子空间 $m$ 的主码字节。中间的 subspace-major 副本使用字节偏移 $mN+p$。最终打包为
\begin{equation}
 \mathsf w[j,p]=\sum_{t=0}^{3}2^{8t}Z[4j+t,p],\qquad
 \operatorname{addr}_{\rm byte}(\mathsf w[j,p])=4(jN+p).
 \label{eq:4-primary-address}
\end{equation}
这里 $j=0,\ldots,M/4-1$。word 布局要求 $M\bmod4=0$，启用该布局的构建对其他值抛出异常；另行编译的字节布局保留原表示。每份打包数组占 $NM$ 字节，但打包期间新旧副本同时存活，占 $2NM$。一遍构建器临时 gather 成 vector-major 的 $[N,M]$ 主码数组，再转置和打包，每份旧副本用完后释放。两遍构建器直接 scatter 到 subspace-major 数组，然后执行相同的最终打包。搜索时的 lane 映射见\Cref{sec:5-3}。

残差码保持 row-major，每向量占 $b_r=\lceil dB_r/8\rceil$ 字节。八比特配置中，行 $p$ 的坐标 $t$ 位于字节偏移 $pb_r+t$。四比特路径中，它位于 $pb_r+\lfloor t/2\rfloor$，偶数坐标在低 nibble，奇数坐标在高 nibble。每物理行另存一个四字节修正项：
\begin{equation}
 \gamma[ p ]=2\langle\widehat p(x_{\rho(p)}),
                              \widehat r(x_{\rho(p)})\rangle.
 \label{eq:4-correction}
\end{equation}
gather 或 scatter 时，其值必须与对应行的编码共同移动。因此，即使两个编码数组采用不同布局，残差访问仍重建主码筛选识别的同一对象。

这一分离适应工作负载的两种访问模式：主码在列表顺序中广泛扫描，完整残差行仅在选定位置读取。设备保留原始 FP32 向量可提供精确重排，但要额外占用 $4Nd$ 字节，并改变所报告的搜索设计。这里精排压缩重建，质量影响由第七章评估。


\section{常驻存储、工作区与索引生命周期}
\label{sec:4-4}
\subsection{存储字节与分配字节}

八比特残差配置具有 int8 量级的常驻表示。主码每子空间一字节时，行载荷为
\begin{equation}
 S_{\rm rows}=N\left(M+\left\lceil\frac{dB_r}{8}\right\rceil+8\right),
 \qquad S_{\rm offsets}=4(n_{\rm list}+1).
 \label{eq:4-row-storage}
\end{equation}
额外八字节保存四字节 ID 与 FP32 修正项。在 $M=d/8$、$B_r=8$ 时，编码占 $9d/8$ 字节，比 SQ8 向量的 $d$ 字节多 12.5\%。因此，无图层次结构是一种设计选择，不能据此声称独有的 int8 压缩优势。

\begin{table}[tbp]
\centering
\caption[单向量存储模型]{每向量字节数。JHQ 含编码、ID 与修正项；SQ8 仅计向量编码；CAGRA int8 模型计向量编码及 32 个四字节邻居。均排除偏移、模型参数、工作区及其他运行时分配，不代表完整进程总量。}
\label{tab:4-bytes}
\footnotesize\setlength{\tabcolsep}{4pt}
\begin{tabular}{@{}rrrrrrrr@{}}
\toprule
$d$ & 主码 & 残差 & ID + 修正 & JHQ & SQ8 & \shortstack{CAGRA\\int8} & FP32 \\
\midrule
768 & 96 & 768 & 8 & 872 & 768 & 896 & 3,072 \\
1,024 & 128 & 1,024 & 8 & 1,160 & 1,024 & 1,152 & 4,096 \\
1,536 & 192 & 1,536 & 8 & 1,736 & 1,536 & 1,664 & 6,144 \\
3,072 & 384 & 3,072 & 8 & 3,464 & 3,072 & 3,200 & 12,288 \\
\bottomrule
\end{tabular}
\end{table}

CAGRA 比较使用图度数 $d_g=32$ 时的 $d+4d_g$ 字节，与测试程序的 32-bit 邻居 ID 一致。在 Stella 上，JHQ 行载荷为 19.20\,GiB，成功的度数 32 CAGRA int8 记录报告索引分配 19,656\,MiB（19.20\,GiB）。尽管计量边界不同，两者接近仍是必要背景。第七章同时报告 int8、fp32 分配与搜索质量及成本。JHQ 的特点是无图构建和查询可调的残差访问预算；降低 $\alpha$ 不释放常驻残差码。

基本 FP32 模型数组中，变换占 $4d^2$ 字节，物化主码码本占 $4MKD_s$，残差码本占 $4MK_r$，IVF 中心与范数占 $4n_{\rm list}(d+1)$。选定路径还分配派生中心副本和编码级值。查询工作区包括两个 $[b_{\rm cap},d]$ FP32 查询矩阵、$[b_{\rm cap},n_{\rm list}]$ 中心分数、LUT、探针及候选缓冲区。可选残差 LUT 另占 $4b_{\rm cap}dK_r$ 字节；默认精排直接计算残差分数。分组工作区遵循\cref{eq:5-group-workspace}，其容量可大于当前扫描所需：已记录的 128 探针样本即使走 query-centric 路径，也保留约 93.8\,MiB。

\begin{table}[tbp]
\centering
\caption[构建与查询对象生命周期]{对象生命周期。大小描述具名数组，不是 CUDA 进程的全部分配；构建阶段见\cref{sec:4-5}。}
\label{tab:4-lifetimes}
\footnotesize\setlength{\tabcolsep}{4pt}
\begin{tabularx}{\linewidth}{@{}>{\raggedright\arraybackslash}p{2.0cm}>{\raggedright\arraybackslash}p{1.7cm}>{\raggedright\arraybackslash}p{1.5cm}*{2}{>{\raggedright\arraybackslash}X}@{}}
\toprule
对象 & 生产者 & 位置 & 生命周期 & 字节／规模 \\
\midrule
输入／样本 & 加载器／测试程序 & 主机 & 构建输入；外部持有 & 原始 $4Nd$；采样子集 \\
变换／模型 & QR／训练 & 主机、GPU & 模型及索引寿命 & 基本 GPU 数组 $4d^2+4MKD_s+4MK_r+4n_{\rm list}(d+1)$ \\
训练缓存 & 训练 & 主机文件 & 跨进程；仍需建库 & 序列化模型状态 \\
输入暂存 & 添加流水线 & Pinned 主机、GPU & add 期间 & 主机 $4n_ba_bd$；GPU 浮点数组 $8n_ba_bd$，另加 scratch \\
ID／偏移 & 排序／lower bound & GPU & 完成后的快照 & $4N+4(n_{\rm list}+1)$ \\
编码／修正项 & 编码／重排 & GPU & 完成后的快照 & $N(M+b_r+4)$ \\
顺序／逆映射 & 排序／求逆 & GPU & 仅构建期间 & 每数组 $4N$，依阶段而定 \\
查询工作区／graph & 搜索／捕获 & GPU、主机 & 复用分配，逐批替换内容 & 依容量而定；graph 引用常驻数组 \\
\bottomrule
\end{tabularx}
\end{table}

\subsection{复用、恢复与失效}

缓存保存旋转、码本和中心。key 混合模型参数及训练样本的等距校验，但不含代码版本，也不覆盖全部数据库字节。因此，复现需要模型来源，并在实现或训练数据变化时显式使缓存失效。搜索实验复用训练状态，冷建记录则构造独立实例；实例间 recall 接近不能证明模型逐字节一致。

缓存命中仍先执行构造器 QR 与库预热，之后训练调用才检查缓存；再加载／上传模型、编码数据并构建列表。Stella 的 6.67\,s add 记录只是该恢复路径的一个组件，不是实测进程重启时间。add 调用包括初始工作区分配；调度校准与搜索预热在这一孤立 add 数字之外。

固定 $k$ 时，改变 $\alpha$ 会改变候选预算，无需重新训练或添加向量。搜索在捕获前扩展不足的工作区，并在实际批大小、候选数、探针数、精排 lane 数、扫描组数或共享界设置变化时重新捕获。校准尝试预算时可绕过捕获。释放索引会销毁 graph 和设备分配；缓冲区由索引持有，在本文评测的顺序调用之间复用。


\section{构建成本与已测试规模}
\label{sec:4-5}
\subsection{有界输入批次与两遍构建}

一次旋转全部 $N$ 个原始向量需要两个 FP32 设备数组，共 $8Nd$ 字节。构建器改用 pinned 主机内存和设备原始／旋转缓冲区分批暂存。若有 $n_b$ 个流水缓冲区、配置 add 行数 $a_{\rm add}$、每数组字节预算
$S_{\rm stage}$，源码暂存行数为
\begin{equation}
 a_{\rm in}=\max\!\left\{1,\min\!\left(a_{\rm add},
 \left\lfloor\frac{S_{\rm stage}}{4d}\right\rfloor\right)\right\},
 \qquad S_{\rm float}=8n_ba_{\rm in}d.
 \label{eq:4-staging-rows}
\end{equation}
默认三个缓冲区、128\,MiB 预算，使输入／旋转设备浮点数组至多占 768\,MiB。
分配乘积和编码 scratch 是独立分配。

一遍路径先编码输入顺序的主码／残差数组，再 gather 成列表顺序，临时与最终压缩副本共存。容量规则预测内存不足时，两遍构建先分配列表、排序 ID、构建偏移，再次读取输入，旋转、编码并直接 scatter 到最终位置。它用额外一次输入扫描和旋转换取较少的同时存活压缩副本。两遍路径的默认分配可在原始输入与反变换中心之间比较；正交性给出相应的实数算术恒等式。

\subsection{阶段占用与进程峰值}

源码的一遍路径分支估计为
\begin{equation}
 H_{\rm single}=N(2b_r+3M+20)+2^{30}\quad\text{bytes}.
 \label{eq:4-peak-heuristic}
\end{equation}
其中依赖 $N$ 的项，是 gather／transpose 时同时存活的十个数组的精确载荷：输入顺序的主码、残差和修正项；分配结果及排序顺序；最终 ID、残差和修正项；vector-major 与 subspace-major 主码副本。共计 $3NM+2Nb_r+20N$ 字节。额外 1\,GiB 是容量启发式的预留，不是其他全部分配的精确核算。

两遍构建在最终 repack 时，同时保留 ID、残差、修正项和新旧主码数组：
\begin{equation}
 S_{\rm repack,2}=N(b_r+2M+8).
 \label{eq:4-two-pass-repack}
\end{equation}
这一表达式对该阶段的这些数组精确成立。更早阶段还保留流水线 scratch，scatter 时另有临时逆置换。因此，repack 表达式本身不是两遍构建的整体峰值。取整个生命周期最大值时，还需计入模型、排序／分配器 scratch、运行时上下文及之后的工作区分配。

\begin{table}[tbp]
\centering
\caption[各构建阶段的行数组显存模型]{行数组的解析占用（GiB），$d=1{,}024$、$M=128$、$B_r=8$。BGE 使用一遍，Stella 使用两遍；Stella 一遍列为反事实容量分析。这些阶段值不含流水线、模型、偏移和运行时成本。}
\label{tab:4-phase-storage}
\small
\begin{tabular}{@{}lrrr@{}}
\toprule
数据集 & 最终行载荷 & 一遍 gather & 两遍 repack \\
\midrule
BGE & 10.90 & 23.05 & 12.11 \\
Stella & 19.20 & 40.59 & 21.32 \\
\bottomrule
\end{tabular}
\end{table}

\Cref{tab:4-phase-storage}量化了压缩副本的节省。Stella 冷日志选择两遍，报告\cref{eq:4-peak-heuristic}为 41.6\,GiB，尾部可用为 30.6\,GiB。另行采样的常驻进程总量包括上下文及查询工作区，BGE 为 11.93\,GiB，Stella 为 20.44\,GiB。容量模型见\cref{sec:4-4}。
补证以20\,ms间隔采样，BGE在add开头观测26.08\,GiB，Stella在repack阶段
观测22.58\,GiB。这是含上下文的设备采样最大值，可能漏掉更短暂分配。

分配乘积容量使用独立配置行数 $a_{\rm assign}$，由
\texttt{JHQ\_ASSIGN\_BATCH} 或默认值选择：列表数至少 8,192 时为 32,768，
否则为 8,192。
\begin{equation}
 S_{\rm dots}=4n_b n_{\rm list}a_{\rm assign}\quad\text{bytes}.
 \label{eq:4-assignment-scratch}
\end{equation}
分配容量不必随 $a_{\rm in}$ 或某次调用的实际处理行数缩小。
\Cref{tab:b-scratch}核对四种默认配置：1,536／3,072 维虽暂存行数不同，
分配乘积均占 3\,GiB。BGE／Stella 在 add 入口附近仍分配三个各 4\,GiB 数组，
共 12\,GiB。
23.05\,GiB的gather载荷属于后续阶段，不能直接再加12\,GiB作为真实峰值。
峰值模型应对各阶段实际存活数组（含流水scratch）取最大值。
Stella两遍repack的较小载荷加模型／上下文，解释了其数据库更大却观测峰值更低。

\subsection{构建测量}

\Cref{sec:7-7}通过重复空缓存构建、较宽启动区间和采样设备峰值评价上述阶段模型，
并报告基线构建与大数据集搜索。计时边界沿用\cref{sec:3-3}；
上述构建算法和分配生命周期决定哪些成本可以重叠。


\section{索引表示与安全契约}
\label{sec:4-6}
历史上，32-bit 的 $n_t d$ 在 $2^{31}$ 处溢出，导致主码尺度归零；后续进程又因 cache key 不含代码版本而加载了错误模型。分块计算范数并以 double 累积修复了溢出。这一故障要求分别核查模型几何、编码行与放置关系，即使粗路由成功也不例外。

\subsection{已完成的表示证据}

10 月 13 日的残差下界 pilot 从 arms6 训练缓存导出 $\Pi$、码本、ID、偏移、修正项及逻辑解码后的打包编码，覆盖完整的 Vogue（932,328 行）、arXiv（2,253,000 行）和 OpenAI-3072（999,000 行）数据库。

码本检查发现可分离性不匹配、标量级值不匹配、残差码本无序对均为零。独立解码导出的主码／残差后，FP64 循环遍历每一行，核对 $\gamma=2\langle\widehat p,\widehat r\rangle$。三个数据库的最大修正项绝对误差为$1.28\times10^{-7}$--$2.77\times10^{-7}$；采样GPU返回top-10分数的最大相对误差为$2.38\times10^{-7}$--$4.85\times10^{-5}$，定义为$|s_{\rm GPU}-s_{64}|/s_{64}$，参考含查询范数常数。这些是导出表示的检查，不是从原始输入向量重新编码；逐数据集诊断保留在仓库补充材料。


在七个工作负载、896 个查询实例中，pilot 对 72,800,590 个路由候选实例检查所提残差下界与 FP64 重建距离，未报告违反；同一集合也支持 FP64 组合分数恒等式检查。GPU 返回分数检查仅覆盖 top-10 槽位，并跳过路由映射中不存在的 ID，因此尚未完整验证输出成员关系。

规范化命名空间并去掉导出方法后，编码、主码／残差训练与构建源码同评测实现一致（\path{scripts/audit_ch4_review.py}）。该对应覆盖共同路径；后续模型身份与大数据集两遍构建仍需单独核查。

\subsection{独立输入与放置审计}
\label{sec:4-representation-audit}

补证导出器在Vogue、arXiv、OpenAI-3072与Stella上使用共同构建路径。
四者全行偏移均单调、端点为0和$N$，ID均构成数据库行的置换。
对存储旋转矩阵做FP64乘法，$\|\Pi\Pi^{\mathsf T}-I\|$最大逐项误差为
$3.96$--$4.34\times10^{-7}$。采样范数检查使用FP64旋转而非GPU GEMM，
因此验证的是存储矩阵的几何性质。\Cref{tab:4-representation-audit}汇总采样重编码。

\par\medskip\noindent\begin{minipage}{\linewidth}
\centering\captionsetup{hypcap=false}
\captionof{table}[原始输入 FP64 重编码审计]{原始输入FP64重编码审计。差异数按坐标决策计数，最后一列为残差不一致
坐标离决策边界的最大距离。结构检查覆盖全行，重编码采用采样。}
\label{tab:4-representation-audit}
\scriptsize\setlength{\tabcolsep}{3pt}
\begin{tabular}{@{}lrrrrr@{}}
\toprule
数据集 & 重编码行数 & 正交误差 & 主码差异 & 残差差异 & 最大边界距离 \\
\midrule
Vogue & 200,000 & 4.11e-07 & 14 & 5,409 & 1.12e-07 \\
arXiv & 100,000 & 4.11e-07 & 1 & 2,515 & 7.31e-08 \\
OpenAI-3072 & 20,000 & 3.96e-07 & 5,745 & 1,646,806 & 2.67e-05 \\
Stella & 100,000 & 4.34e-07 & 16 & 3,064 & 1.03e-07 \\
\bottomrule
\end{tabular}

\end{minipage}\par\medskip

768／1,024维上的差异均在决策边界$1.12\times10^{-7}$内。
OpenAI-3072则明显不同：残差不一致占所查坐标的2.68\%，最远离边界
$2.67\times10^{-5}$，与TF32旋转一致，因此不能将所有数据集称为FP32等价。
残差比较以已存主码位为条件，避免重复计入主码翻转。

三个较小数据库在同一模型下强制一遍／两遍构建，按ID关联的主码和残差字节完全
相同，correction差至多$8.94\times10^{-8}$；列表分配分别有1、1、40行不同。
Stella只验证自然两遍路径。导出器不含IVF质心，因此列表成员决策本身仍未独立验证，
结构正确不等于已证明分配决策正确。

\subsection{契约覆盖与剩余检查}
% Keep this float after the preceding fixed-position recoding-audit table.
\suppressfloats[t]

\begin{table}[tbp]
\centering
\caption[表示契约与审计覆盖]{纳入pilot及补证输入审计后的表示契约，区分存储几何、采样算术与行放置的覆盖。}
\label{tab:4-contracts}
\footnotesize\setlength{\tabcolsep}{4pt}
\begin{tabularx}{\linewidth}{@{}>{\raggedright\arraybackslash}p{1.8cm}*{2}{>{\raggedright\arraybackslash}X}@{}}
\toprule
边界 & 所需性质 & 证据／剩余检查 \\
\midrule
C0：几何 & 共同变换，浮点等距误差有界 & 四数据集已查存储矩阵正交性；区分TF32 GPU算术 \\
C1：码本／修正 & 笛卡尔级值、有序残差码本，修正项匹配解码编码 & 三个数据库 pilot 已完成；全行修正误差见\cref{sec:4-6} \\
C1：输入编码 & 存储编码对应预期原始输入行 & 四份原始输入采样已做FP64重编码；量化边界与TF32差异已量化 \\
C1：放置 & 偏移覆盖行，ID 为置换，字段遵循同一置换 & 全行偏移／ID核对及三份按ID关联的路径比较；质心成员决策未验证 \\
执行 & 索引、同步和分数顺序契约有效 & 已有查询与 sanitizer 检查，有限覆盖见\cref{sec:7-4} \\
\bottomrule
\end{tabularx}
\end{table}

构建要求 $M\mid d$、$D_s\mid B$、$M\bmod4=0$，LUT 验证使用实际码字切片。
输入须有限、计数可表示并预留哨兵位置。ID／位置用 signed 32-bit；较大计数与地址乘积用更宽算术。

审计使用导出模型而非按 seed 重生成。T4 检查计算分数选择，与输入重编码分开
（\cref{sec:7-strict-selection}）；质心成员关系、空列表与对抗尾批次仍待验证。

